%% Generated by make_tex.py from SMOTE25_Method_Map_1.html (28 Oct 2026). Do not edit by hand.
%% Each account was written from a full reading of the PDF on 28 Oct 2026 (AI reader); check against the paper.
\subsection{Binary only (24 Methods)}
\label{sec:methods-B}
\begin{description}
\item[SMOTE (2002)] \parencite{chawla2002smote} This paper defines the reference rule and changes none of the six decision points: every minority sample is a seed with N/100 points each (a random N\% of seeds when N $<$ 100), the partner is drawn from the k = 5 minority nearest neighbours, and $\lambda \sim$ U(0,1) on the segment (Sec 4.2, p.328; algorithm, p.329). Generation runs once before training, combined with random undersampling of the majority (Sec 4.3, p.331). SMOTE-NC changes only the distance used to find neighbours and sets nominal values by a vote of the neighbours (Sec 6.1, p.348).
  \textit{Steps changed:} none (the reference rule).
  \classev{9 datasets, all used as two classes (Sec 5.1, pp.333-334; Table 2, p.334): Pima, Phoneme, Adult, E-state, Satimage, Forest Cover, Oil, Mammography, Can; Adult also run with SMOTE-NC (Sec 6.1, pp.348-349). The paper focuses on two-class problems; a multiclass use is stated but not run (pp.333-334). \textit{Reduction:} Satimage (6 classes): smallest class vs the other 5 merged (Sec 5.1 item 5, p.333). Forest Cover (7 classes): only Ponderosa Pine and Cottonwood/Willow kept, the other classes dropped (Sec 5.1 item 6, pp.333-334).}
\item[SMOTE-Tomek, SMOTE-ENN (2004)] \pcite{batista2004study} SMOTE first balances the training set (Sec. 5, p.25); then Tomek-link pairs are removed from both classes, or, in the ENN variant, every example misclassified by its 3 nearest neighbours is removed from both classes (Sec. 4.3, pp.23-24; Fig. 2). Seeds, neighbours and placement are as in SMOTE; the only addition is cleaning after generation.
  \textit{Steps changed:} Post-filtering.
  \classev{15 two-class datasets built from 13 UCI sources (Sec. 5, Table 3, pp.24-25); the paper considers only two-class problems (Sec. 4.1, p.22). C4.5, 10-fold CV, AUC (Tables 4-7). \textit{Reduction:} Smallest class vs the rest merged (Sec. 5, p.24), e.g. Vehicle van, New-thyroid hypo, E.Coli iMU, Satimage 4, Flag white, Glass, Nursery; Letter gives Letter-a (a vs rest) and Letter-vowel (all vowels vs rest); Splice gives two one-vs-rest sets (Table 3).}
\item[Borderline-SMOTE (2005)] \pcite{han2005borderline} For each minority point its m nearest neighbours in the whole training set are found (Sec. 3, Steps 1-2, p.882); points with only majority neighbours are treated as noise, and points with at least half majority neighbours form the DANGER set, which are the only seeds. In v1 each DANGER seed interpolates towards s of its k = 5 minority neighbours with $\lambda \sim$ U(0,1) as in SMOTE (Step 4, p.882). v2 also interpolates towards the seed's nearest majority neighbour with $\lambda$ in (0, 0.5), so those points stay near the seed (pp.882-883).
  \textit{Steps changed:} Seed selection, Neighbour choice, Placement rule.
  \classev{4 datasets used as two-class problems (Sec. 4, Table 2, p.883): Circle (simulated), Pima, Satimage (class 4 vs the rest), Haberman. The paper studies only two-class problems (Sec. 4, p.883). C4.5, 3 x 10-fold CV, TP rate and F-value (p.884). \textit{Reduction:} Satimage: class 4 vs all remaining classes merged (Sec. 4, p.883; Table 2).}
\item[Cluster-SMOTE (2006)] \pcite{cieslak2006cluster} k-means is run on the minority examples only and SMOTE is applied separately inside each cluster, after which original and synthetic minority points are merged (Sec. III-C, p.4). The interpolation partner therefore always comes from the seed's own cluster; seeds, the SMOTE rate per cluster and uniform placement on the segment are as in SMOTE (Sec. III-B, p.3; Sec. V-A, p.5).
  \textit{Steps changed:} Neighbour choice.
  \classev{Two datasets built by the authors from Notre Dame campus traffic (Sec. II, p.2): Packets (346,620 instances, 2,106 alerts) and Destinations (4,080 instances, 147 alerts) (Table I, p.2), both treated as intrusive vs not intrusive (Sec. II, p.2). Cluster-SMOTE was run only on Packets with 3 and 5 k-means clusters (Sec. V-A, Fig. 4, p.5); on Destinations the minority formed one cluster (Sec. V-B, pp.5-6). \textit{Reduction:} No benchmark multiclass dataset; SNORT attack-type labels are collapsed to intrusive vs not (Sec. II, p.2).}
\item[ADASYN (2008)] \pcite{he2008adasyn} ADASYN keeps the SMOTE generator: a partner is drawn from the seed's K nearest minority neighbours and the point is seed + $\lambda$(partner - seed), with K = 5 in the experiments (Eq. 5, p.1324; Sec. III-C, p.1325). What changes is allocation: the total G = (ml - ms) beta is shared out as $g_{i}$ proportional to the share of majority points among each seed's K neighbours (Eqs. 2-4, pp.1323-1324), instead of an equal number per seed.
  \textit{Steps changed:} Sample allocation.
  \classev{Algorithm defined for labels y in \{1,-1\} (Sec. II, p.1323). Artificial 2-class Gaussian set (Fig. 1, p.1324) and 5 UCI datasets used as two-class problems (Table I, Sec. III-A, pp.1324-1325): Vehicle, Pima, Vowel, Ionosphere, Abalone. A multiclass extension is described (Sec. III-D, p.1327) but never run. \textit{Reduction:} Vehicle (4 classes): van vs the rest (p.1324). Vowel (11 classes): first vowel vs the rest (p.1325). Abalone (29 classes): only class 18 vs class 9 kept (p.1325). Pima and Ionosphere are natively two-class.}
\item[LS (2008)] \pcite{cieslak2008ls} LS leaves the SMOTE generator unchanged but applies it per segment: a Hellinger-distance tree of height at most 2 splits the training data into up to 4 segments (Alg 1, Sec 3, pp4-5), and a 5-fold CV wrapper searches each segment for its own undersampling level (0-100\%) and SMOTE level (0-500\%) (Alg 2, Sec 4, p5). The number of synthetic points per seed therefore varies by segment, and generation runs inside a wrapper search rather than once before training. Seeds, partner and placement inside a segment are as in SMOTE.
  \textit{Steps changed:} Sample allocation, Inside a pipeline.
  \classev{20 datasets (Table 2, p7): real-world from Chawla et al. 2002, LIBSVM, the KDD'99 Intrusion subset and biological sets, plus UCI (Sec 5.3, pp7-8). All used as two-class problems with one minority percentage each (Table 2); Intrusion keeps u2r vs normal (Sec 5.3, p8). Motivating experiments use artificial two-class radial data (Sec 2, Fig 2-3, Table 1, pp2-4); Pima is the worked example (Fig 4, p6). \textit{Reduction:} Multiclass UCI sources reduced to smallest class vs the rest; Letter uses vowels vs consonants (Sec 5.3, p8). Which datasets were converted is not stated.}
\item[Safe-Level-SMOTE (2009)] \pcite{bunkhumpornpat2009safelevel} Each minority point has a safe level, the number of minority points among its k nearest neighbours (Eq 7, p478), and the ratio of the seed's and partner's safe levels bounds the gap (Fig 1, Sec 3, pp478-479): ratio 1 gives U(0,1) as in SMOTE, above 1 pulls the point towards the seed, below 1 towards the partner. When the partner's safe level is 0 the seed is copied, and when both are 0 nothing is generated, so noise pairs are skipped. The partner is still drawn uniformly from the 5 nearest minority neighbours and nothing is removed afterwards.
  \textit{Steps changed:} Seed selection, Placement rule.
  \classev{Studies only the two-class case (Sec 1, p475). Two UCI datasets (Table 3, p480): Satimage (6,435; 626 positive) and Haberman (306; 81 positive). k=5, 10-fold CV, C4.5, Naive Bayes and SVM (Sec 4, p480). \textit{Reduction:} Satimage: class 4 vs all other classes merged (Sec 4, p480).}
\item[RAMOBoost (2010)] \pcite{chen2010ramoboost} The generator is SMOTE's segment interpolation towards one of the seed's k2 nearest minority neighbours (Eq. 1, p.1626; Algorithm 1 step 5). In each boosting round, seeds are drawn with replacement from a ranking that weights each minority point by a logistic function of the number of majority points among its k1 neighbours (steps 2-4, Eqs. 2-3), so border points receive more synthetic samples. Generation is repeated in every AdaBoost.M2 round on a minority subset resampled by the boosting weights (steps 1, 6-10).
  \textit{Steps changed:} Seed selection, Sample allocation, Inside a pipeline.
  \classev{19 UCI/ELENA datasets, all run as two-class problems (Tables I-II, pp.1630-1631); only two-class problems considered (Sec. IV, p.1630); multiclass is future work (Sec. V, p.1641). \textit{Reduction:} One-vs-rest (e.g. Wine 1 vs 2+3, Vehicle van, Segment brickface, Letter Z vs A-Y), merged classes (Texture, Page\_blocks) or two classes kept and the rest dropped (Abalone 18 vs 9, Yeast POX vs CYT) (Table II, p.1631).}
\item[INOS (2013)] \pcite{cao2013inos} INOS splits the synthetic quota between two generators: 70\% from ESPO, a regularised Gaussian of the minority covariance that does not interpolate (Sec. 3.2.2, Eq. 13; r = 0.7, Sec. 4.3), and 30\% from ADASYN in a signal space with the null space removed (Sec. 3.3, Eqs. 14-15). The ADASYN part allocates points per seed by the share of negatives among its neighbours and interpolates uniformly towards one of the 5 nearest positives, as in SMOTE. Afterwards, Tomek-like noise links between synthetic positives and nearest negatives are removed (Sec. 3.4, pp.5-6).
  \textit{Steps changed:} Sample allocation, Post-filtering.
  \classev{Restricted to two-class time series (p.2; Sec. 3, p.3). 48 two-class apportions from 7 UCR datasets (Adiac, FaceAll, 50words, SLeaf, Wafer, Yoga, TwoPats), each a separate binary task (Sec. 4.2, pp.7-8, Table 1). \textit{Reduction:} Adiac (37 classes), FaceAll (14), 50words (50), SLeaf (15): one class positive, the rest merged as negative, one apportion per chosen class (Sec. 4.2, p.8).}
\item[MWMOTE (2014)] \pcite{barua2014mwmote} Seeds are limited to an informative minority set near the borderline majority points, after dropping minority points whose neighbours are all majority (Alg. 1 steps 1-6; Sec. 4.1, p.5). Each seed gets a weight from its closeness to the majority class and its density, and seeds are drawn in proportion to it, so counts per seed vary (Alg. 1 steps 7-9, 12a; Sec. 4.2, pp.6-7). The partner is a random member of the seed's minority cluster from average-linkage clustering, not one of its 5 nearest neighbours (Alg. 1 steps 10, 12b; Secs. 4.3-4.4). Placement is uniform on the segment and nothing is removed afterwards.
  \textit{Steps changed:} Seed selection, Sample allocation, Neighbour choice.
  \classev{4 artificial two-class domains (Sec. 5.3, pp.11-12) and 20 UCI datasets used as two-class problems (Sec. 5.4, p.13, Table 5); networks with 2 output neurons (Sec. 5.2). Multiclass use named as future work (Sec. 6, p.19). \textit{Reduction:} Multiclass sources converted to two-class problems with the chosen minority and majority classes listed in Table 5 (Sec. 5.4, p.13); the table could not be extracted from the text.}
\item[RBBoost (2015)] \pcite{diezpastor2015randombalance} Each ensemble member gets a training set of the original size with a random class split (Alg. 1 lines 6-7, p.98): the class that must grow is filled with SMOTE (k = 5; Sec. 3, p.97) and the other is randomly undersampled, and when the split is inverted SMOTE grows the original majority class (Alg. 1 lines 12-15; Sec. 4.1, p.98). Partner and placement are as in SMOTE. Generation is redone for every member of a plain or bagging ensemble, or in every AdaBoost.M2 round (RB-Boost, Alg. 2, Sec. 4.2, p.100).
  \textit{Steps changed:} Seed selection, Sample allocation, Inside a pipeline.
  \classev{Stated for two-class problems (Abstract; Sec. 1, p.97); Algorithms 1-2 use labels y in \{-1,+1\} (pp.98, 100). Two-class 2-D Gaussian simulation (Sec. 5, p.100). 86 binary datasets: 20 HDDT (Table 2, p.102) and 66 KEEL (Table 3, p.103), 5x2-fold CV, J48 (Sec. 6, p.101). Multiclass extension listed as future work only (Sec. 7, p.110). \textit{Reduction:} 66 KEEL sets derived from 16 sources (footnote 7, p.101) as one-vs-rest or merged-group splits (e.g. ecoli-0-1-4-7\_vs\_2-3-5-6, yeast-2\_vs\_4, page-blocks-1-3\_vs\_4); HDDT sets such as letter, pendigits, satimage are presumably binarised, reduction not described.}
\item[SMOTE-IPF (2015)] \pcite{saez2015smoteipf} Generation is plain SMOTE with k = 5, oversampling to a 50/50 balance (Sec. 3.2, p.189). The only change is a filter applied afterwards: the Iterative-Partitioning Filter, an ensemble of C4.5 models on 9 partitions with majority voting, repeated until fewer than 1\% of examples are removed in 3 rounds, which removes noisy original and synthetic examples of both classes (Sec. 3.1, pp.188-189; Sec. 3.3, p.190).
  \textit{Steps changed:} Post-filtering.
  \classev{30 synthetic two-class 2-D datasets (Sec. 4.1, pp.190-191) and 9 real datasets, each one minority class vs the rest, 8 from KEEL (Sec. 4.2, p.191; Table 1, p.192), with class or attribute noise added. \textit{Reduction:} Multiclass sources turned into two classes by joining one or more classes as positive and the rest as negative (Sec. 4.2, p.191), likely ecoli, newthyroid and cleveland.}
\item[A-SUWO (2016)] \pcite{nekooeimehr2016asuwo} Minority points are grouped by a complete-linkage clustering that will not merge two minority sub-clusters across a majority sub-cluster (Sec. 3.1, pp.26-28). Each sub-cluster's quota is proportional to its cross-validated classification error (Sec. 3.2, Eqs. 3.2-3.3), seeds are drawn with weights favouring closeness to the majority class (Sec. 3.3, Eqs. 3.4-3.9), and the partner is one of the seed's nearest neighbours in the same sub-cluster (Alg. 3.1, p.33). Placement stays uniform on the segment (Eq. 3.10).
  \textit{Steps changed:} Seed selection, Sample allocation, Neighbour choice.
  \classev{Drive PDF is the 2016 dissertation, Chapter 3 being the A-SUWO paper. Defined for binary problems (p.24; Sec. 3.5, p.46). 16 datasets, each a minority class vs the rest (Table 3.1, p.35); multiclass use only future work (Sec. 3.5, p.46). \textit{Reduction:} 12 multiclass sources reduced to named class(es) vs all other classes (e.g. Vehicle van, Ecoli pp, Glass 1, Segment WINDOW) (Table 3.1, p.35).}
\item[CGMOS (2016)] \pcite{zhang2016cgmos} CGMOS keeps SMOTE's interpolation step unchanged (Sec. 2.3, p.1625). Seeds are drawn in proportion to a weight: the average relative change in kernel-density posterior certainty over all training points if a minority point were added at that seed (Eqs. 5-6, p.1625). Sec. 3 proves that the expected certainty gain is at least that of SMOTE on the training data (Theorem 1).
  \textit{Steps changed:} Seed selection.
  \classev{Set up for two classes only (Sec. 2, p.1624; Sec. 5, p.1629). 30 UCI datasets used as two-class problems (Sec. 4.1, p.1626; Table 1, p.1627). \textit{Reduction:} Multiclass datasets keep the smallest class as minority and merge the rest (Sec. 4.1, p.1626); which datasets is not stated.}
\item[k-means SMOTE (2018)] \pcite{douzas2018kmeans} k-means clusters the whole input space, and only clusters whose majority-to-minority ratio is below a threshold are kept, so seeds are minority points inside minority-dominated clusters (Sec. 3.1, PDF pp.7-8; Algorithm 1). Samples are shared across kept clusters in proportion to their sparsity, so sparse clusters get more (Sec. 3.1, PDF pp.8-9). Inside each cluster plain SMOTE runs on that cluster's points only, so partners come from the same cluster; placement is unchanged and there is no filtering. With one cluster and no threshold the method reduces to SMOTE (Sec. 3.2, PDF p.9).
  \textit{Steps changed:} Seed selection, Sample allocation, Neighbour choice.
  \classev{All experiments two-class; a multiclass extension is said to be possible but not covered (Sec. 1, PDF p.4). 90 datasets (Abstract): 12 UCI sets plus randomly undersampled copies, 19 KEEL two-class sets and 10 simulated MADELON sets (Sec. 4.4, PDF p.14; Table 2, PDF p.15), plus 5 two-class 2-D toy sets (Figs. 10-14). \textit{Reduction:} Multiclass UCI sources (e.g. breast tissue, ecoli, glass, iris, libra, segment, vehicle, wine) turned into one-vs-rest problems with the smallest class as minority (Sec. 4.4, PDF p.14); KEEL sets are ready-made binary splits (Table 2).}
\item[G-SMOTE (2019)] \pcite{douzas2019gsmote} Seeds and allocation are SMOTE-like (Sec. 4.1, PDF p.9). The surface point that sets the radius is chosen by a selection strategy: a random minority neighbour as in SMOTE, the nearest majority point, or whichever of the two is closer (Algorithm 1, PDF p.8; Sec. 4.3 step 4). The segment is replaced by a point drawn in a hypersphere around the seed with that radius, truncated and deformed into a hyper-spheroid along the seed-to-surface axis (Sec. 4.2; Sec. 4.3 steps 6-9, PDF pp.14-16). With full truncation and deformation and the minority strategy the region becomes SMOTE's segment (Sec. 5.4, PDF p.24).
  \textit{Steps changed:} Neighbour choice, Placement rule.
  \classev{All experiments two-class: 28 base datasets (13 UCI, 13 KEEL, 2 simulated MADELON) plus undersampled copies, 69 in total (Sec. 5.1, PDF p.18; Table 1, pp.18-19); binary AUC, F-score and G-mean (Sec. 5.2). \textit{Reduction:} Multiclass sources (e.g. iris, wine, glass, vehicle, segmentation, ecoli, libras, vowel, led, dermatology) appear as two-class sets in Table 1; how they were reduced is not stated.}
\item[AB-SMOTE (2020)] \pcite{almajzoub2020absmote} As in Borderline-SMOTE, only danger minority points (at least half but not all of their M nearest neighbours are majority) are seeds (Sec. II.B, p.32; Sec. III, pp.32-33). The difference is that the k nearest neighbours, and so the partner, are searched only within the borderline set (Sec. III, Figs. 5-6, p.33). Placement is P'i + gap (P'j - P'i) with gap in [0,1], as in SMOTE (Sec. III, p.33).
  \textit{Steps changed:} Seed selection, Neighbour choice.
  \classev{Framed for binary classification (title; Sec. I, p.31). Eight two-class datasets (Table I, p.33): Abalone9-18, Flare, Wine Quality (KEEL); Haberman (UCI); Customer Churn Wireless Telecom (CrowdAnalytix); Telco Customer Churn, Employee Attrition, Sales Win/Loss (IBM). J48, 5-fold CV (Sec. IV, pp.33-34). \textit{Reduction:} No reduction described; Abalone9-18, Flare and Wine Quality are KEEL two-class versions of multiclass sources.}
\item[HDE (2021)] \pcite{chen2021hde} The partner is one of the k = 5 nearest majority instances (nearest enemies), found with a mix of Euclidean distance and how often earlier ensemble members disagree on the two points (Sec. 3.3, Eqs. 3-4, pp.6-7). The step is x + r($x_{i}$ - x) with r drawn from N(0,1), so points cluster near the seed and can fall on either side of it (Eq. 5, p.8). Generation runs inside every member of a sequential bagging ensemble (Secs. 3.1-3.4).
  \textit{Steps changed:} Neighbour choice, Placement rule, Inside a pipeline.
  \classev{Targets binary tasks (Abstract; Sec. 2; Conclusion). 42 binary sets: 28 KEEL and 14 LibSVM (Sec. 4.1, Table 1, pp.9-10). \textit{Reduction:} KEEL x\_vs\_y and one-vs-rest splits of multiclass sources (Table 1); reduction not described.}
\item[FW-SMOTE (2022)] \pcite{maldonado2022fwsmote} The only change is how the k nearest minority neighbours are found: Euclidean distance is replaced by a feature-weighted Minkowski distance whose weights come from a supervised feature ranking, and low-ranked features can be dropped from the distance (Eqs. 4-5, 9-12, pp.4-5; Algorithm 1 steps 2-9). Seeds, the number per seed, the random partner and placement on the segment are as in SMOTE (Algorithm 1 steps 10-14), and the method reduces to SMOTE with equal weights and p = 2 (p.6).
  \textit{Steps changed:} Neighbour choice.
  \classev{Defined for labels \{-1,+1\} (Sec. 3.3, p.5). 42 two-class datasets: 30 low-dimensional UCI/KEEL and 12 high-dimensional microarray sets (Table 1, p.7; Sec. 4.1, pp.6-7); binary AUC and G-mean. \textit{Reduction:} Low-dimensional sets are mostly KEEL 'x vs y' and one-vs-rest versions of multiclass sources (Table 1, p.7); the paper does not describe the reduction.}
\item[KNNOR (2022)] \pcite{islam2022knnor} Seeds are only the densest d\% of minority points, ranked by distance to their k-th minority neighbour, with d chosen automatically (Alg. 1, Sec. 3.1, pp.3-4). Instead of one random partner, the point moves step by step through all k nearest minority neighbours, each step bounded by the smallest majority-minority distance (Eq. 5, Alg. 2, Sec. 3.2, pp.6-8). A candidate is kept only if all its k nearest neighbours are minority points (Alg. 3, Sec. 3.3, p.8).
  \textit{Steps changed:} Seed selection, Neighbour choice, Placement rule, Post-filtering.
  \classev{19 two-class datasets (18 KEEL plus CM1; Table 2, Sec. 4.1, p.11) and two-class image demos (UTKFace; CIFAR-10 airplane vs rest; Sec. 6, p.14). Multiclass use only future work (Sec. 7, p.16). \textit{Reduction:} KEEL splits of multiclass UCI sources (Sec. 4.1, p.11); CIFAR-10 one-vs-rest.}
\item[SVDD-WSMOTE (2022)] \pcite{tao2022svddw} Noisy or overlapped minority points inside the majority SVDD sphere are removed before generation (Secs. 4.1-4.2). Each remaining seed gets a number of points proportional to a weight combining closeness to the boundary and sparsity (Eqs. 21-24), and the partner is drawn from its 5 nearest neighbours in proportion to their weights (Eq. 25, p.25). Placement is SMOTE's; synthetic points that fall inside the SVDD boundary are rejected and regenerated, and any still missing are filled by copying the seed (Alg. steps 36-52, pp.26-27).
  \textit{Steps changed:} Seed selection, Sample allocation, Neighbour choice, Post-filtering.
  \classev{Four 2-D simulated two-class sets (Sec. 5.1), 24 binary datasets from 16 UCI sources (Table 2, pp.34-35) and bearing-fault data with fault types against normal (Table 8, pp.44-45). Defined for one majority and one minority class (Eqs. 7-8, 24); multiclass is future work (Sec. 6). \textit{Reduction:} One-vs-others or class-vs-class (Abalone 16 vs 6, Yeast3-5) (Sec. 5.3; Table 2).}
\item[CISO (2026)] \pcite{yadav2025ciso} Anchors are drawn with replacement in proportion to a safety score, so each anchor gets a varying number of points (Eqs. 8-9, p.8). Two different minority neighbours of the anchor are drawn with weights from their clearance and a q-Gaussian locality kernel (Eqs. 10-11). The new point lies on the segment between those two neighbours, not between anchor and neighbour, with $\lambda$ drawn from a truncated q-Gaussian instead of U(0,1) (Eqs. 12-14, p.12). Nothing is removed afterwards. Running time is 39 times SMOTE's at the median (Sec. 6.7, p.29).
  \textit{Steps changed:} Seed selection, Sample allocation, Neighbour choice, Placement rule.
  \classev{Binary labels only (Sec. 3.1, p.5). 37 binary KEEL datasets for evaluation, 5 KEEL datasets for tuning and 8 larger UCI datasets (Sec. 5.1, p.15; Tables 1-2). Multiclass is future work (Sec. 7, p.34). \textit{Reduction:} Multiclass sources made binary as smallest class vs rest (Sec. 5.1; Table 1 caption), e.g. Waveform, Statlog Landsat, Page Blocks, Letter; KEEL sets are ready-made binary versions.}
\item[Loyalty-SMOTE (2026)] \pcite{hu2026loyalty} Generation is plain SMOTE (k = 5) applied only to boundary minority samples whose Loyalty, an inverse-distance-weighted balance of same-class and other-class neighbours, lies between 0 and m (Sec. 3.3). Samples with negative Loyalty are treated as possible noise and kept out of generation, so the partner pool is also restricted. Afterwards Loyalty is recomputed: flagged minority samples still below 0 are removed and flagged majority samples are removed by an LOF test (Sec. 3.4, Eqs. 10-12).
  \textit{Steps changed:} Seed selection, Neighbour choice, Post-filtering.
  \classev{30 binary UCI/KEEL datasets (group 1, Table 2) plus other two-class groups, and 5 multiclass UCI datasets (group 4, Table 5: yeast 10 classes, HR 3, employee 3, Drug200 4, MEDV 9) with results in Sec. 4.4.4 and Tables 15-18. \textit{Reduction:} Group 1 includes KEEL x\_vs\_y splits and one-vs-rest versions of multiclass sources (Table 2); reduction not described. \textit{Several classes:} On multiclass data the largest class is the majority and the smallest the minority (Sec. 4.4.4); SMOTE oversamples only the smallest class (Sec. 3.3); low-Loyalty samples are relabelled to the class with the highest Attraction before generation (Eq. 9, Sec. 3.2). Whether the middle classes stay in the data is read differently by two checks; only one class is ever oversampled, and results use binary metrics.}
  \pnote{Provisional: the PDF check read this as Binary + Multiclass; coded Binary only under rule 1 (only the smallest class is ever oversampled, binary metrics). Decision for P. Category decision needed: only the smallest class is oversampled on multiclass data, scored with binary metrics.}
\item[QC-SMOTE (2026)] \pcite{upman2026qcsmote} Seeds are drawn in proportion to a trust score combining safe level, minority density and distance from the majority class (Eqs. 2-6, Sec. 3.2, p.5), and the number of attempts per seed also follows trust (Eq. 7, Sec. 3.3, p.6). Each candidate is still $x_{i}$ + $\lambda$($x_{j}$ - $x_{i}$) with a 5-NN minority partner, but $\lambda \sim$ U(0, $\alpha_{r}$) shrinks with overlap and the best of $K_{r}$ candidates is kept by a purity and clearance score (Eqs. 8-18, Secs. 3.4-3.5, pp.6-7). A candidate that fails the thresholds is replaced by a copy of the most trusted point near the seed (Eqs. 19-21, Sec. 3.6, p.7).
  \textit{Steps changed:} Seed selection, Sample allocation, Placement rule, Post-filtering.
  \classev{30 binary datasets from the VS-SMOTE benchmark (KEEL-style), IR 1.82 to 32.73 (Table 1, p.8; Sec. 4.1, p.8). CatBoost, repeated stratified CV, oversampling on training folds only (Secs. 4.2-4.3, p.8). A one-vs-rest use is described (Sec. 3.1, p.4; Sec. 3.7, p.7) but never run. \textit{Reduction:} Uses ready-made KEEL two-class splits (one-vs-rest such as ecoli1-4, glass0/1/6, vehicle0-3, yeast1/3/4/5; merged-class splits such as glass-0-1-2-3\_vs\_4-5-6, yeast-1-2-8-9\_vs\_7) (Table 1, p.8); the paper does not describe the reduction.}
\end{description}
\subsection{Multiclass only (5 Methods)}
\label{sec:methods-M}
\begin{description}
\item[SCUT (2015)] \pcite{agrawal2015scut} Synthetic points come from standard SMOTE run inside each under-represented class (Fig. 1); the only change is how many points each class gets, enough to reach the mean class size m with a rate computed per class (Sec. 2). Seeds, partners within the class and placement are as in SMOTE, and generation is a one-off step before training. The EM cluster-based undersampling acts on real points of the large classes, not on synthetic ones.
  \textit{Steps changed:} Sample allocation.
  \classev{8 multiclass datasets with all classes kept: 7 KEEL (Thyroid 3 classes, Lymphography 4, Pageblocks 5, Dermatology 6, Autos 6, Ecoli 8, Yeast 10) and UCI Wine Quality (7) (Sec. 3.1, Table 1). No two-class experiment. 10-fold CV with J48, SMO, Naive Bayes, IBk (Sec. 3). \textit{Several classes:} Mean class size m is the common target: classes below m are SMOTEd up to m with a per-class rate, classes above m are reduced to m by EM clustering and near-uniform sampling across clusters (Sec. 2, Fig. 1; Sec. 3.1). No decomposition, no ensemble.}
\item[SMOM (2017)] \pcite{zhu2017smom} SMOM keeps SMOTE-style interpolation between a seed and one of its k1 = 12 neighbours (Supp. Secs. 3, 5). For trapped instances, each neighbour direction gets a selection weight that penalises directions whose segment sphere contains instances of other classes (Algorithm S1; Supp. Sec. 3), so directions are drawn by weight rather than uniformly. NBDOS clustering decides which instances are outstanding (uniform directions) or trapped (weighted directions).
  \textit{Steps changed:} Neighbour choice.
  \classev{Read from the supplement only: the main PDF in Drive contains the INOS paper, not SMOM. 27 multiclass datasets, UCI and ordinal-regression benchmarks, all classes kept, with per-group minority and majority metrics (Supp. Tables S3-S10, S17-S20). No two-class experiment appears in the supplement. \textit{Several classes:} Each minority class is processed in turn against all other classes and oversampled until all classes have the same size (Supp. Sec. 7, Table S21); neighbours are searched over all classes; no decomposition or ensemble.}
  \pnote{Main PDF mislinked in Drive; placement law and allocation not confirmed from the main text.}
\item[EMDO (2021)] \pcite{yao2021emdo} Each minority class is covered by Mahalanobis ellipsoids and its quota is split across them by inverse density times purity (Eqs. 32-34, p.11). In pure ellipsoids points are drawn along the ellipsoid axes without interpolation (Eqs. 37-40). In impure ellipsoids a modified Borderline-SMOTE interpolates: only borderline points are seeds, and the partner comes from their 5 Mahalanobis neighbours within the ellipsoid (Eqs. 41-42, p.12).
  \textit{Steps changed:} Seed selection, Sample allocation, Neighbour choice, Placement rule.
  \classev{9 KEEL/UCI sets with all classes kept (3 to 6 classes; Table 1, p.13) and 2 sensor sets (Shuttle 7, Mafalda 3; Table 6, p.16). Only two-class data is a toy illustration (Fig. 3). In 5 of the 11 sets only one class is oversampled. \textit{Several classes:} Any class whose size ratio to the largest class is below the threshold is grown to N\_max/IR (Eq. 1, p.4); each minority class gets its own ellipsoids fitted against all other classes (Eqs. 15-18).}
\item[GLOS (2023)] \pcite{han2023glos} Only one of GLOS's three generators interpolates: AMBO, used for borderline seeds (Alg. 2, p.2076). Its partner is any instance inside a radius-R sphere, possibly of another class, and the step is capped at the seed's nearest-neighbour distance rather than covering the whole segment (Alg. 2 line 12). The number of points per seed comes from discreteness and is split across classes by their discreteness ratios. The other two generators place points on Mahalanobis contours (MMO) or at random within a radius (RRO).
  \textit{Steps changed:} Seed selection, Sample allocation, Neighbour choice, Placement rule.
  \classev{21 multiclass UCI/KEEL datasets with 3 to 11 classes kept (Sec. 4.1, Table 1). One parameter-tuning run on Sonar (Fig. 7, p.2079) also reports GLOS results; the paper does not state Sonar's class count (two classes in the standard UCI version). \textit{Several classes:} Classes are labelled minority or majority by comparing their discreteness with the dataset-wide value (Eqs. 1, 2, 5); minority seeds are typed safe, borderline or small-disjunct and each type gets its own generator (MMO, AMBO, RRO) (Algs. 1-3).}
  \pnote{Provisional: Multiclass only; would be Binary + Multiclass if the two-class Sonar tuning run (Fig. 7) counts. Decision for P.}
\item[MKC-SMOTE (2024)] \pcite{wang2024mkc} A minority point is a seed only if at least half of its 5 nearest neighbours, taken from all classes, share its class (Sec. III step 2; Alg. 1). Instead of a random minority partner, the method uses the centroid of the seed's neighbours and places new points on the line from that centroid through the seed, possibly beyond the seed (Sec. III steps 3-5; Alg. 1 lines 23-26). It generates twice the needed number and keeps the candidates nearest the centroid (Sec. III step 6, p.5).
  \textit{Steps changed:} Seed selection, Neighbour choice, Placement rule, Post-filtering.
  \classev{21 UCI datasets, all multiclass (Abstract; Sec. IV-A-1, p.5; Table 1), all classes kept and scored with multiclass MG and MAUC (Eqs. 6-7, p.6); 5-fold CV, CART and SVM. No two-class experiment. \textit{Several classes:} Every class with N\_max/N\_i $>$ 1.45 is oversampled towards the largest class size, each class independently, with the neighbour search over all classes (Alg. 1; Sec. III, pp.4-5). No decomposition.}
\end{description}
\subsection{Binary + Multiclass (8 Methods)}
\label{sec:methods-BM}
\begin{description}
\item[SMOTEBoost (2003)] \pcite{chawla2003smoteboost} The generator is unchanged SMOTE (Sec 2, p.109), but it runs inside every AdaBoost.M2 round and the synthetic points are discarded after the round (Sec 3, p.110; Fig. 1, p.111). New points are created by sampling seeds from the round's boosting distribution, which the authors say focuses generation on misclassified or boundary minority cases (pp.110, 116). Allocation, partner, placement and the absence of filtering are as in SMOTE.
  \textit{Steps changed:} Seed selection, Inside a pipeline.
  \classev{4 datasets (Sec 4, Table 2, p.112): KDDCup-99 with all 5 classes kept (DoS, Probe, Normal, U2R, R2L); Mammography and Phoneme (two-class); Satimage (reduced to two). Results tabulated only for U2R (Table 3, p.113); R2L omitted for space. RIPPER base learner (p.111). \textit{Reduction:} Satimage: smallest class vs the rest merged (p.112). KDDCup-99 not reduced; only DoS and Normal subsampled for size (p.112). \textit{Several classes:} AdaBoost.M2 core over labels \{1..C\} (Fig. 1, p.111); on KDDCup-99 SMOTE is applied separately to each minority class (U2R at 100/300/500\%, R2L at 100\%) (p.112); the other three classes are not oversampled.}
\item[SSRBF, DSRBF (2011)] \pcite{fernandeznavarro2011dsrbf} The generator is unchanged SMOTE: every original pattern of the chosen class is a seed joined to one of its k = 5 nearest same-class neighbours (Sec. 1, p.1821; Sec. 3.4, p.1826). What changes is which class is oversampled and how many times, set by class size (SSRBF) or by the lowest per-class sensitivity of the current best network (DSRBF). In DSRBF generation runs inside the training loop; SSRBF runs before training.
  \textit{Steps changed:} Sample allocation, Inside a pipeline.
  \classev{13 datasets (Sec. 4.1, p.1827; Table 3, p.1828): 5 two-class (Hepatitis, BreastC, Haberman, CYTvsPOX, ME2vsOther) and 8 with all classes kept (Newthyroid 3, Balance 3, Saureus4 4, Anneal 5, Glass 6, Zoo 7, E. coli 8, Yeast 10). SSRBF and DSRBF run on all 13 (Table 4) with one softmax RBF network, no decomposition. \textit{Reduction:} CYTvsPOX and ME2vsOther are derived from Yeast (Sec. 4.1, p.1827). \textit{Several classes:} One class is oversampled at a time: SSRBF enlarges the currently smallest class J times before training; DSRBF pauses the memetic training at set generations and oversamples the class with the lowest sensitivity, each time by its original size (Secs. 3.3-3.4, pp.1825-1827).}
\item[GAMO (2019)] \pcite{mullick2019gamo} GAMO creates each sample as a convex combination of all training points of a minority class (Eq. 1, Sec. 3.2, p.1697), not as a point on a seed-neighbour segment, so there is no single seed and no k-NN partner: a generator network learns the weights adversarially to place points near class peripheries that the classifier gets wrong. Generation runs inside end-to-end training that alternates generator, classifier and discriminator updates (Alg. 1, p.1698), usually in a learned feature space.
  \textit{Steps changed:} Seed selection, Sample allocation, Neighbour choice, Placement rule, Inside a pipeline.
  \classev{Main paper: seven image datasets, all multiclass with selected classes kept: MNIST and Fashion-MNIST (10 classes, IR 100; Sec. 4.1, Table 1), CIFAR10 and SVHN (10 classes; Sec. 4.2, Table 2), CelebA and LSUN (5 classes, small and large; Sec. 4.3, Table 3), SUN397 (50 classes; Sec. 4.4, Table 4). Supplement (Sec. 4, Tables 7-8, p.8): GAMO also run on three two-class sets (Abalone19, IJCNN1, Magic) and two multiclass sets with all classes kept (Chess 18 classes, Cover Type 7). \textit{Reduction:} None in the main paper. Supplement uses Abalone19, the KEEL two-class version of Abalone (reduction not described); IJCNN1 and Magic are two-class by nature. \textit{Several classes:} The c-1 minority classes are each filled up to the majority count (Sec. 3.1, p.1697); a shared class-conditional unit feeds one instance-generation unit per minority class, which outputs convex weights over that class's instances (Sec. 3.2, Eq. 1); labels of generated points are drawn in proportion to (Pc - Pi) (Alg. 1). One c-class classifier, no decomposition.}
\item[RSMOTE (2021)] \pcite{chen2021rsmote} Minority points whose k neighbours are all majority are dropped as noise, and the remaining seeds are split into borderline and safe groups by 2-means on their relative density (Defs. 1-3, pp.401-402; Algorithm 1, p.403). Each group gets a share matching its size, and within a group each seed gets a weight (k-m)/(k+1), so safer seeds produce more points (Step 3, p.402). Points are made with SMOTE's own routine, uniform on the segment, once before training (Algorithm 1 line 22).
  \textit{Steps changed:} Seed selection, Sample allocation.
  \classev{Two-class: 4 synthetic 2-D sets (Sec. 4.3.1, p.406) and 27 UCI sets used as binary (Table 2, p.409), each at 0-30\% label noise. Multiclass: 8 UCI sets with all classes kept, 3 to 5 classes (Table 7, p.416; Sec. 4.3.3, pp.413-414), macro-averaged metrics (Eqs. 7-10, p.405). Same generator in both. \textit{Reduction:} Several of the 27 binary sets come from multiclass UCI sources (e.g. ecoli, iris, wine, segment, vehicle; Table 2, p.409); the reduction is not described but counts suggest one class vs the rest. \textit{Several classes:} Every class smaller than the largest is oversampled with RSMOTE up to the largest class size (Sec. 4.3.3, p.413); classification uses diversified one-against-one classifiers (p.414).}
\item[NROMM (2023)] \pcite{liu2023nromm} Clustering splits minority samples into inland, borderline and trapped (Sec. 3.2). Inland and borderline seeds are interpolated, inland ones with a partner from the same cluster and borderline ones with a nearby majority neighbour; trapped seeds get points from an optimised embedding, not interpolation (Sec. 3.3, Eqs. 2-4). The number of points per seed follows (1 - local density) times the class-size gap, categorical values come from a VDM-constrained set near the seed (Sec. 4.1), and borderline and trapped majority samples are removed afterwards.
  \textit{Steps changed:} Seed selection, Sample allocation, Neighbour choice, Placement rule, Post-filtering.
  \classev{20 KEEL multiclass datasets with 3 to 10 classes kept (Sec. 5.1; Table 2, p.7) and 10 two-class KEEL datasets (D1-D10) in the validity check (Sec. 5.3, p.9; Table 6; Supp. Table 1). \textit{Reduction:} D1-D10 are KEEL two-class versions of multiclass sources (one-vs-rest such as Yeast4, Ecoli4; merged splits such as Ecoli-3-5 vs 0-6-7) (Supp. Table 1). \textit{Several classes:} One-versus-ensemble: classes sorted by size, each non-largest class oversampled against the original and synthetic data of all larger classes (Sec. 4.2.1; Fig. 4); afterwards inland samples of the overlap region are removed (Sec. 4.2.2).}
  \pnote{Core vs Partial debatable: trapped seeds (6-84\% per dataset) are not interpolated. The binary and multiclass variants appear to be run on separate data.}
\item[OREM, OREM-M (2023)] \pcite{zhu2023orem} OREM keeps every minority point as a seed with equal counts (Alg. 1). It replaces the k-NN partner set with 'assistant seeds' from a candidate region that grows until q = 5 consecutive majority points are met, keeping a candidate only if the sphere on the seed-candidate segment holds no majority point, so borderline majority points can be partners (Alg. 1; Sec. 3.1.2). The point is $x_{i}$ + g($x_{s}$ - $x_{i}$) with g $\sim$ U[0,1], halved when the partner is a majority point. OREM-M adds class-quota sampling and a nearest-neighbour rejection filter (Alg. 2), and OREMBoost runs the generator inside each AdaBoost.M2 round (Alg. 3).
  \textit{Steps changed:} Seed selection, Sample allocation, Neighbour choice, Placement rule, Post-filtering, Inside a pipeline.
  \classev{OREM on 28 two-class UCI datasets (Sec. 4.1.1, Table 2, p.6174); OREM-M on 21 multiclass datasets with 3 to 10 classes kept (Sec. 4.2.1, Table 6, p.6177); OREMBoost on both (Sec. 4.3, Tables 10-11). \textit{Reduction:} 20 of OREM's 28 sets come from multiclass sources reduced to one or merged classes vs the rest (Table 2). \textit{Several classes:} OREM-M (Alg. 2, Sec. 3.2, p.6172): only classes with imbalance degree above 1.2 are oversampled; one point at a time, a class is drawn in proportion to its remaining quota, and a synthetic point is rejected if its nearest neighbour belongs to another minority class. Native multiclass, not OvO/OvR.}
  \pnote{Variants: OREM is B, OREM-M is M, OREMBoost is BM; the paper is BM through OREMBoost and the union.}
\item[MC-MBRC (2024)] \pcite{ma2024mcmbrc} Seeds are only the safe samples of each class, after membership-based noise removal (Alg. 1; Sec. 3). The partner is one of the 5 same-class neighbours among the safe and border samples, and placement is SMOTE's $x_{i}$ + U(0,1)($x_{j}$ - $x_{i}$) (Alg. 2). Each seed gets a random number of points up to the class quota divided by the number of safe seeds. Afterwards, a CCR-style energy sphere moves overlapping and generated samples of larger classes towards their class centre (Algs. 3-4, p.6).
  \textit{Steps changed:} Seed selection, Sample allocation, Neighbour choice, Post-filtering.
  \classev{20 multiclass UCI/KEEL datasets with 3 to 10 classes kept (Sec. 4.2, Table 2, p.7); 60 synthetic two-class datasets (Sec. 5.2, Fig. 5) and 50 more in the ablation (Sec. 5.4, Fig. 7). \textit{Several classes:} Membership to each class centre splits every class into safe, overlapping and noisy samples (Alg. 1); every class gets a quota N\_max minus its overlap count (Alg. 2); cleaning goes through classes by descending size, treating all larger classes as majority (Alg. 4). No decomposition.}
\item[GDHS (2025)] \pcite{yan2025gdhs} Generation interpolates between a minority seed and one of its 5 same-class neighbours (Eq. 3, p.3), but each seed-direction pair gets a weight that favours safe seeds (same-class share among neighbours, Eq. 2) and directions whose synthetic region holds few other-class points (Eqs. 5-10). Each minority class is filled to the mean class size by weighted random draws rather than a fixed number per seed (Alg. 1). Afterwards, original majority points in the synthetic regions, or overlapping other majority classes, are removed (Sec. 3.2, Eqs. 11-16).
  \textit{Steps changed:} Seed selection, Sample allocation, Neighbour choice, Post-filtering.
  \classev{20 datasets (Sec. 4.1, Table 1, pp.6-7): 14 KEEL multiclass sets with all classes kept (3 to 10 classes, e.g. car, glass, yeast, zoo) and 6 two-class software-defect sets (camel-1.4, camel-1.6, hive-0.12.0, velocity-1.5, xalan-2.4, xerces-1.2). All three variants run on all 20 sets (Tables 5-8, pp.10-11). \textit{Reduction:} None; the defect sets are natively two-class. \textit{Several classes:} Mean class size m is the target (Eq. 1): classes below m are oversampled up to m, classes above m are cleaned by one of three strategies (Sec. 3.2, pp.5-6). No decomposition.}
\end{description}
